% GENERATED FILE. Edit canonical lessons, then run npm run generate:books.
% canonical-source-hash: fnv1a64:5da55a82f6dc57d6
% canonical-lessons: TE-C480-muncettu, TE-C480-pongiporlu, TE-C480-tegipovu, TE-C480-kurukupovu, TE-C480-irukkupovu

\chapter{Floodwater}
\label{ch:te-floodwater}
\hlchaptermodality{480}

\begin{chapteropening}
\textbf{By the end of this chapter:} \emph{I can say to inundate, to overflow, to snap, to get bogged down and to get stuck.}

Complete the last lesson of chapter 480: I can say to inundate, to overflow, to snap, to get bogged down and to get stuck.
\end{chapteropening}

\section[\texorpdfstring{\te{ముంచెత్తు} (muñcettu)}{ముంచెత్తు (muñcettu)}]{\te{ముంచెత్తు} (muñcettu) — to inundate; to flood}

\label{lesson:TE-C480-muncettu}



\begin{quote}
Before the new one: say the Telugu for to grow dim, then the Telugu for spray.
\end{quote}

\subsection*{You'll want to know: \te{ముంచెత్తు}}
\textbf{\te{ముంచెత్తు}} — \emph{muñcettu} — \textquotedblleft{}to inundate; to flood\textquotedblright{}.

\begin{quote}
\te{వరద} \te{నీరు} \te{గ్రామాన్ని} \te{ముంచెత్తింది}.
\end{quote}

\begin{quote}
\emph{varada nīru grāmānni muñcettindi.}
\end{quote}

\begin{quote}
The floodwater inundated the village.
\end{quote}

Past (I …): \te{ముంచెత్తాను}, \emph{muñcettānu}.

\te{ముంచు}, to dip, + \te{ఎత్తు}, to lift: to plunge under. Also said of a crowd or of gifts: \te{పూలతో} \te{ముంచెత్తారు}.

\begin{cousinweb}
\textbf{\te{ముంచెత్తు}} (\emph{muñcettu}) is native Dravidian: \te{ముంచు}, \textquotedblleft{}to plunge, to immerse\textquotedblright{} (DEDR 4993), and \te{ఎత్తు}, \textquotedblleft{}to lift, to raise\textquotedblright{}.
\end{cousinweb}

\begin{grammarlens}[title={What you've built}]
One more word for getting around.
\end{grammarlens}

\subsection*{Your turn}
\begin{itemize}
\raggedright
  \item \emph{Say it:} \emph{muñcettu}
  \item \emph{Say it:} \emph{muñcettu}, once more
  \item \emph{Say it:} \emph{varada nīru grāmānni muñcettindi.} — The floodwater inundated the village.
  \item \emph{Recall:} say the Telugu for to grow dim, then the Telugu for spray, then say \emph{muñcettu} again
\end{itemize}

\begin{tcolorbox}[breakable,colback=teal!4,colframe=teal!35!black,title={Before you move on}]
What does \te{ముంచెత్తు} mean? (\textquotedblleft{}To inundate; to flood\textquotedblright{}.) Say it once more.
\end{tcolorbox}
\section[\texorpdfstring{\te{పొంగిపొర్లు} (poṅgiporlu)}{పొంగిపొర్లు (poṅgiporlu)}]{\te{పొంగిపొర్లు} (poṅgiporlu) — to overflow}

\label{lesson:TE-C480-pongiporlu}



\begin{quote}
Before the new one: say the Telugu for spray, then the Telugu for to inundate.
\end{quote}

\subsection*{You'll want to know: \te{పొంగిపొర్లు}}
\textbf{\te{పొంగిపొర్లు}} — \emph{poṅgiporlu} — \textquotedblleft{}to overflow\textquotedblright{}.

\begin{quote}
\te{భారీ} \te{వర్షాలకు} \te{వాగులు} \te{పొంగిపొర్లుతున్నాయి}.
\end{quote}

\begin{quote}
\emph{bhārī varṣālaku vāgulu poṅgiporlutunnāyi.}
\end{quote}

\begin{quote}
With the heavy rains the streams are overflowing.
\end{quote}

Past (it …): \te{పొంగిపొర్లింది}, \emph{poṅgiporlindi}.

\te{పొంగి}, swelling up, + \te{పొర్లు}, to roll over: water rising over its banks. Its subject is a thing, so the past shown is the third person: \te{పొంగిపొర్లింది}, it overflowed.

\begin{cousinweb}
\textbf{\te{పొంగిపొర్లు}} (\emph{poṅgiporlu}) is native Dravidian: \te{పొంగు}, \textquotedblleft{}to bubble up, to boil over\textquotedblright{} (DEDR 4469), and \te{పొర్లు}, \textquotedblleft{}to roll over\textquotedblright{}, the twin of \te{దొర్లు}.
\end{cousinweb}

\begin{grammarlens}[title={What you've built}]
One more word for getting around.
\end{grammarlens}

\subsection*{Your turn}
\begin{itemize}
\raggedright
  \item \emph{Say it:} \emph{poṅgiporlu}
  \item \emph{Say it:} \emph{poṅgiporlu}, once more
  \item \emph{Say it:} \emph{bhārī varṣālaku vāgulu poṅgiporlutunnāyi.} — With the heavy rains the streams are overflowing.
  \item \emph{Recall:} say the Telugu for spray, then the Telugu for to inundate, then say \emph{poṅgiporlu} again
\end{itemize}

\begin{tcolorbox}[breakable,colback=teal!4,colframe=teal!35!black,title={Before you move on}]
What does \te{పొంగిపొర్లు} mean? (\textquotedblleft{}To overflow\textquotedblright{}.) Say it once more.
\end{tcolorbox}
\section[\texorpdfstring{\te{తెగిపోవు} (tegipōvu)}{తెగిపోవు (tegipōvu)}]{\te{తెగిపోవు} (tegipōvu) — to snap; to be cut off}

\label{lesson:TE-C480-tegipovu}



\begin{quote}
Before the new one: say the Telugu for to inundate, then the Telugu for to overflow.
\end{quote}

\subsection*{You'll want to know: \te{తెగిపోవు}}
\textbf{\te{తెగిపోవు}} — \emph{tegipōvu} — \textquotedblleft{}to snap; to be cut off\textquotedblright{}.

\begin{quote}
\te{వరదలో} \te{రోడ్డు} \te{తెగిపోయింది}.
\end{quote}

\begin{quote}
\emph{varadalō rōḍḍu tegipōyindi.}
\end{quote}

\begin{quote}
The road was cut off in the flood.
\end{quote}

Past (it …): \te{తెగిపోయింది}, \emph{tegipōyindi}.

Said of a rope, a wire, a road or a bond that breaks apart. Its subject is a thing, so the past shown is the third person: \te{తెగిపోయింది}, it snapped.

\begin{cousinweb}
\textbf{\te{తెగిపోవు}} (\emph{tegipōvu}) is a Telugu compound: \te{తెగు}, \textquotedblleft{}to break off, to be cut\textquotedblright{}, whose origin is not settled, and the native \te{పోవు}, \textquotedblleft{}to go\textquotedblright{} (DEDR 4572).
\end{cousinweb}

\begin{grammarlens}[title={What you've built}]
One more word for getting around.
\end{grammarlens}

\subsection*{Your turn}
\begin{itemize}
\raggedright
  \item \emph{Say it:} \emph{tegipōvu}
  \item \emph{Say it:} \emph{tegipōvu}, once more
  \item \emph{Say it:} \emph{varadalō rōḍḍu tegipōyindi.} — The road was cut off in the flood.
  \item \emph{Recall:} say the Telugu for to inundate, then the Telugu for to overflow, then say \emph{tegipōvu} again
\end{itemize}

\begin{tcolorbox}[breakable,colback=teal!4,colframe=teal!35!black,title={Before you move on}]
What does \te{తెగిపోవు} mean? (\textquotedblleft{}To snap; to be cut off\textquotedblright{}.) Say it once more.
\end{tcolorbox}
\section[\texorpdfstring{\te{కూరుకుపోవు} (kūrukupōvu)}{కూరుకుపోవు (kūrukupōvu)}]{\te{కూరుకుపోవు} (kūrukupōvu) — to get bogged down; to sink in}

\label{lesson:TE-C480-kurukupovu}



\begin{quote}
Before the new one: say the Telugu for to overflow, then the Telugu for to snap.
\end{quote}

\subsection*{You'll want to know: \te{కూరుకుపోవు}}
\textbf{\te{కూరుకుపోవు}} — \emph{kūrukupōvu} — \textquotedblleft{}to get bogged down; to sink in\textquotedblright{}.

\begin{quote}
\te{బురదలో} \te{లారీ} \te{చక్రాలు} \te{కూరుకుపోయాయి}.
\end{quote}

\begin{quote}
\emph{buradalō lārī cakrālu kūrukupōyāyi.}
\end{quote}

\begin{quote}
The lorry wheels sank into the mud.
\end{quote}

Past (I …): \te{కూరుకుపోయాను}, \emph{kūrukupōyānu}.

Said of wheels in mud, and of a person stuck in debt or work.

\begin{cousinweb}
\textbf{\te{కూరుకుపోవు}} (\emph{kūrukupōvu}) is a Telugu compound whose first part has an origin that is not settled; the second is the native \te{పోవు}, \textquotedblleft{}to go\textquotedblright{} (DEDR 4572).
\end{cousinweb}

\begin{grammarlens}[title={What you've built}]
One more word for getting around.
\end{grammarlens}

\subsection*{Your turn}
\begin{itemize}
\raggedright
  \item \emph{Say it:} \emph{kūrukupōvu}
  \item \emph{Say it:} \emph{kūrukupōvu}, once more
  \item \emph{Say it:} \emph{buradalō lārī cakrālu kūrukupōyāyi.} — The lorry wheels sank into the mud.
  \item \emph{Recall:} say the Telugu for to overflow, then the Telugu for to snap, then say \emph{kūrukupōvu} again
\end{itemize}

\begin{tcolorbox}[breakable,colback=teal!4,colframe=teal!35!black,title={Before you move on}]
What does \te{కూరుకుపోవు} mean? (\textquotedblleft{}To get bogged down; to sink in\textquotedblright{}.) Say it once more.
\end{tcolorbox}
\section[\texorpdfstring{\te{ఇరుక్కుపోవు} (irukkupōvu)}{ఇరుక్కుపోవు (irukkupōvu)}]{\te{ఇరుక్కుపోవు} (irukkupōvu) — to get stuck}

\label{lesson:TE-C480-irukkupovu}



\begin{quote}
Before the new one: say the Telugu for to snap, then the Telugu for to get bogged down.
\end{quote}

\subsection*{You'll want to know: \te{ఇరుక్కుపోవు}}
\textbf{\te{ఇరుక్కుపోవు}} — \emph{irukkupōvu} — \textquotedblleft{}to get stuck\textquotedblright{}.

\begin{quote}
\te{వరదలో} \te{చాలా} \te{కార్లు} \te{ఇరుక్కుపోయాయి}.
\end{quote}

\begin{quote}
\emph{varadalō cālā kārlu irukkupōyāyi.}
\end{quote}

\begin{quote}
Many cars got stuck in the flood.
\end{quote}

Past (I …): \te{ఇరుక్కుపోయాను}, \emph{irukkupōyānu}.

\te{ఇరుకు}, narrow, tight, + -\te{కొని} + \te{పోవు}: caught in a tight place.

\begin{cousinweb}
\textbf{\te{ఇరుక్కుపోవు}} (\emph{irukkupōvu}) is native Dravidian: DEDR 524 lists Telugu \emph{iṟuku}, \textquotedblleft{}to be squeezed, to get stuck in\textquotedblright{}, here with \te{పోవు}, \textquotedblleft{}to go\textquotedblright{} (DEDR 4572).
\end{cousinweb}

\begin{grammarlens}[title={What you've built}]
One more word for getting around.
\end{grammarlens}

\subsection*{Your turn}
\begin{itemize}
\raggedright
  \item \emph{Say it:} \emph{irukkupōvu}
  \item \emph{Say it:} \emph{irukkupōvu}, once more
  \item \emph{Say it:} \emph{varadalō cālā kārlu irukkupōyāyi.} — Many cars got stuck in the flood.
  \item \emph{Recall:} say the Telugu for to snap, then the Telugu for to get bogged down, then say \emph{irukkupōvu} again
\end{itemize}

\begin{tcolorbox}[breakable,colback=teal!4,colframe=teal!35!black,title={Before you move on}]
What does \te{ఇరుక్కుపోవు} mean? (\textquotedblleft{}To get stuck\textquotedblright{}.) Say it once more.
\end{tcolorbox}
